\begin{examplebox}
	\textbf{\textcolor{CExample}{例 1（基础）}}

	\textbf{\textcolor{CExample}{题目：}}
	平行四边形 $ABCD$ 中，$AB=3$，$AD=5$，$\angle BAD=60^\circ$，求 $AC^2+BD^2$。

	\textbf{\textcolor{CExample}{解题步骤：}}
	\begin{description}[style=nextline, leftmargin=2.8em, labelsep=0.5em, itemsep=0.45em]
		\item[\textbf{第一步（识别条件）：}] 四边形 $ABCD$ 为平行四边形，已知邻边长 $3,5$，夹角 $60^\circ$。
			\par\textbf{理由：} 由题目直接给出。
		\item[\textbf{第二步（套用恒等式）：}] 由 $AC^2+BD^2 = 2(AB^2+AD^2)$。
			\[
				AC^2+BD^2 = 2(AB^2+AD^2)
			\]
			\par\textbf{理由：} 平行四边形对角线平方和恒等式，与夹角无关。
		\item[\textbf{第三步（计算数值）：}] $AB^2=9$，$AD^2=25$，代入得：
			\[
				\begin{aligned}
					AC^2+BD^2 &= 2(9+25)\\
					&= 68
			\end{aligned}
			\]
			\par\textbf{理由：} 算术运算。
	\end{description}

	\textbf{\textcolor{CExample}{关键结论：}}
	\textbf{$AC^2+BD^2 = 68$}

	\medskip
	\textbf{\textcolor{CExample}{例 2（稍变形）}}

	\textbf{\textcolor{CExample}{题目：}}
	平行四边形 $ABCD$ 中，$AB=4$，$AC=6$，$BD=2\sqrt{5}$，求边 $AD$ 的长。

	\textbf{\textcolor{CExample}{解题步骤：}}
	\begin{description}[style=nextline, leftmargin=2.8em, labelsep=0.5em, itemsep=0.45em]
		\item[\textbf{第一步（建立方程）：}] 利用恒等式 $AC^2+BD^2 = 2(AB^2+AD^2)$。
			\[
				AC^2+BD^2 = 2(AB^2+AD^2)
			\]
			\par\textbf{理由：} 恒等式将已知量与未知量关联。
		\item[\textbf{第二步（代入数据）：}] $6^2+(2\sqrt{5})^2 = 2(4^2+AD^2)$，即 $36+20=2(16+AD^2)$。
			\[
				36+20 = 2(16+AD^2)
			\]
			\par\textbf{理由：} 平方运算。
		\item[\textbf{第三步（解出 $AD$）：}] $56 = 2(16+AD^2)$，所以 $28=16+AD^2$，$AD^2=12$。
			\[
				\begin{aligned}
					AD^2 &= 12,\\
					AD &= 2\sqrt{3}
			\end{aligned}
			\]
			\par\textbf{理由：} 等式两边同除、移项、开方（边长取正）。
	\end{description}

	\textbf{\textcolor{CExample}{关键结论：}}
	\textbf{$AD = 2\sqrt{3}$}

	\medskip
	\textbf{\textcolor{CExample}{例 3（综合应用）}}

	\textbf{\textcolor{CExample}{题目：}}
	平行四边形 $ABCD$ 中，$AB=3$，$AD=5$，$AC=7$，求 $BD$ 的长及 $\cos\angle BAD$。

	\textbf{\textcolor{CExample}{解题步骤：}}
	\begin{description}[style=nextline, leftmargin=2.8em, labelsep=0.5em, itemsep=0.45em]
		\item[\textbf{第一步（求 $BD$ 平方）：}] 由恒等式 $AC^2+BD^2 = 2(AB^2+AD^2)$。
			\[
				\begin{aligned}
					7^2+BD^2 &= 2(3^2+5^2)\\
					49+BD^2 &= 68\\
					BD^2 &= 19
			\end{aligned}
			\]
			\par\textbf{理由：} 代入已知，求解。
		\item[\textbf{第二步（余弦定理）：}] 在 $\triangle ABD$ 中，应用余弦定理。
			\[
				\begin{aligned}
					BD^2 &= AB^2+AD^2 - 2\cdot AB\cdot AD\cdot\cos\angle BAD\\
					19 &= 34 - 30\cos\angle BAD\\
					30\cos\angle BAD &= 15\\
					\cos\angle BAD &= \frac12
			\end{aligned}
			\]
			\par\textbf{理由：} 解方程求余弦值。
		\item[\textbf{第三步（结果）：}] 由 $\cos\angle BAD = \frac12$ 得 $\angle BAD = 60^\circ$，且 $BD=\sqrt{19}$。
			\[
				\cos\angle BAD = \frac12
			\]
			\[
				\angle BAD = 60^\circ
			\]
			\par\textbf{理由：} 特殊角识别。
	\end{description}

	\textbf{\textcolor{CExample}{关键结论：}}
	\textbf{$BD=\sqrt{19}$，$\cos\angle BAD = \frac12$，$\angle BAD=60^\circ$}
\end{examplebox}
